\documentclass[conference]{IEEEtran}
\usepackage[utf8]{inputenc}
\usepackage[T1]{fontenc}
\usepackage{cite}
\usepackage{amsmath,amssymb}
\usepackage{graphicx}
\usepackage{booktabs}
\usepackage{url}

\title{Verifier‑Grounded Generate‑Verify‑Repair for Intent\ensuremath{\rightarrow}Configuration NetOps Tasks: A Paired Mininet Emulation Study}

\author{
\IEEEauthorblockN{Autonomous AI Researcher}
\IEEEauthorblockA{CNSM 2026 Agentic AI Researcher Track}
}

\begin{document}

\maketitle

\begin{abstract}
We preregistered and executed a provenance‑traced paired confirmatory study (N = 200 intents; 400 paired episodes) to evaluate whether a deterministic verifier‑grounded generate--verify--repair (GVR) loop reduces operationally detectable policy‑violation errors when a hosted LLM produces device configuration sequences from operator intent in a Mininet emulation. Execution used shared\_initial\_candidate semantics (one sampled initial generation per intent reused by both baseline and guarded). Baseline success (no detected policy violation) = 196/200 (98\%); guarded success = 200/200 (100\%). Paired risk difference (guarded - baseline) = 0.02; preregistered exact McNemar two‑sided p = 0.125 (primary confirmatory test). An exploratory paired bootstrap (seed = 1729, 10,000 resamples) yields a 95\% percentile CI = [0.005, 0.04]. Four verifier‑triggered repairs occurred (total hosted model calls = 204), giving mean extra model calls per intent = 0.02 and mean extra calls per avoided violation = 1.0. All reported numeric values, provider‑call accounting, validator traces, and analysis artifacts are drawn verbatim from the archived execution and analysis bundle. We interpret the preregistered confirmatory result conservatively and discuss construct, internal, and external validity limits and operational implications.
\end{abstract}

\section{Introduction}
Automating intent\ensuremath{\rightarrow}configuration translation is an active NetOps research direction because natural‑language or intent descriptions are a natural operator interface but can be transformed incorrectly by large language models (LLMs). Mechanically checkable faults introduced by LLMs---syntactic errors, topology/reachability violations, and ACL/policy mistakes---can lead to availability loss or exposure if applied to devices without validation [1], [4]. Generate--Verify--Repair (GVR) patterns pair an LLM generator with a deterministic verifier: the verifier checks mechanizable constraints and, if violations occur, issues localized diagnostics used to request corrective edits from the model [2], [3]. Prototype work across code and domain‑specific program generation suggests verifier‑grounded repair reduces mechanically verifiable errors; however, questions remain about scaling such pipelines to heterogeneous multi‑vendor template families, the concrete hosted‑LLM cost per avoided violation, and how to interpret paired comparisons when sampling semantics reuse the initial model candidate.

Research question and contribution. Our preregistered research question is: Does a verifier‑grounded generate--verify‑repair loop significantly reduce operationally detectable policy‑violation errors produced by a hosted LLM when generating network configurations for operator intents across heterogeneous multi‑vendor template families and topology patterns in a Mininet‑based emulation? Contributions: (1) a provenance‑traced confirmatory paired experiment (C1‑GVR‑multivendor‑2026‑v1) with shared\_initial\_candidate semantics executed on N = 200 intents (400 episodes); (2) audited provider and verifier accounting showing repair costs and concrete hosted‑model calls; (3) artifact‑grounded statistical inference (exact McNemar test) and exploratory bootstrap interval (seed = 1729, 10,000 resamples); (4) a focused analysis of failure modes, threats to validity, and operationally grounded recommendations consistent with archived artifacts.

\section{Related Work}
Recent literature on LLMs for NetOps and AIOps converges on two observations that are directly relevant to verifier‑grounded pipelines. First, surveys and domain reviews document rapid uptake of LLMs for network tasks while highlighting recurring gaps in evaluation practice: missing provenance, limited verifier integration, and scarce multi‑vendor or emulation‑scale evidence [5], [6]. These reviews synthesize heterogeneous reports (research prototypes, bench demonstrations, and applied surveys) and conclude that while LLMs reduce human effort for drafting configurations, the community lacks standardized, provenance‑traced benchmarks that pair generators with deterministically auditable verifiers. This limits cross‑study comparability and weakens claims about operational safety or generalization across vendors and topologies [5], [6].

Second, targeted configuration studies and benchmarks show feasibility but also expose common failure modes. Prior NetOps evaluations (for example, work investigating LLM requirements for correct router configuration synthesis and the NetConfEval benchmark) document frequent classes of mechanically checkable failures---syntax/grammar mistakes, incorrect sequencing that violates workflow invariants, and topology/reachability errors when configurations contradict the network state [1], [4]. These studies typically evaluate generator accuracy on curated tasks and emphasize the need for automated checking; however, many prototype benchmarks either omit a closed‑loop repair stage or do not provide provenance traces linking initial outputs, verifier diagnostics, and any applied repairs. In practice, reported improvements in such studies are therefore difficult to attribute to deterministic validator interventions versus differences in prompting or sampling.

A cross‑domain body of work motivates verifier‑grounded mitigation architectures (generate--verify--repair). Recent methodological proposals and prototype demonstrations in code and domain‑specific program generation argue that a verifier capable of expressing mechanizable constraints plus localized diagnostic feedback enables targeted repairs that reduce mechanically detectable errors [2], [3]. These works provide conceptual and single‑case empirical support: they show that when a verifier can (i) deterministically check the artifact against explicit constraints and (ii) return actionable, localized diagnostics, a bounded repair policy can converge failing samples to valid artifacts. Importantly, these prior results also identify the core dependency of GVR effectiveness on verifier coverage: if a verifier cannot encode an error class, that class remains a residual risk even with repeated repair attempts [2], [3].

How our study relates. Building on the empirical observations and methodological proposals above, our experiment addresses a concrete gap: scalable, provenance‑traced paired evaluations that measure the marginal effect of deterministic verification and bounded repair across heterogeneous template families and topology patterns. Unlike prior NetOps benchmarks that either evaluate single‑pass generation or lack explicit provenance between generator outputs and verifier actions [1], [4], or like GVR proposals that focus on single‑domain demonstrations [2], [3], our paired design (with shared\_initial\_candidate semantics) isolates the verifier/repair contribution by reusing the identical initial LLM sample for both conditions and then measuring the final guarded outcome after permissive, verifier‑triggered repair. This combination---multi‑vendor synthetic templates exercised in Mininet, deterministic verifier traces, provable provider‑call accounting, and preregistered paired inference---directly responds to the literature's call for provenance and auditable assessments of verifier‑assisted generation [5], [6]. At the same time, we adopt the caution emphasized in the surveyed and prototype literature: verifier expressivity bounds the classes of errors that can be detected and repaired, and mechanical verification cannot by itself guarantee absence of higher‑order semantic misinterpretations outside the verifier's rule set [2], [3], [5].

\section{Methodology}
Preregistration and formal estimand. We preregistered the study as C1‑GVR‑multivendor‑2026‑v1 before any model calls (preregistration SHA‑256: 9e0fc992\allowbreak{}66583597\allowbreak{}5b01289f\allowbreak{}7841926e\allowbreak{}f6abf358\allowbreak{}4252c5d0\allowbreak{}459ba0ec\allowbreak{}f166bb50). The primary estimand is the paired\_success\_rate\_difference\_guarded\_minus\_baseline: the per‑intent difference in the probability that an intent produces zero operationally detectable policy violations under guarded (GVR) relative to baseline (LLM‑only). The confirmatory test declared in the preregistration is the exact McNemar two‑sided test at \ensuremath{\alpha} = 0.05 applied to the paired binary success indicators.

Sampling, stratification, and fixed sampling semantics. Task selection was deterministic and stratified across three synthetic vendor‑template families and four topology patterns (12 strata). N = 200 unique operator intents were fixed by the deterministic task generator before any model calls (task\_index\_profile selection is recorded in the execution manifest). The preregistration and adapter execution contract explicitly used shared\_initial\_candidate semantics: exactly one sampled initial generation per intent was produced and this identical initial candidate was reused by both baseline and guarded episodes. Under these semantics the baseline outcome equals the deterministic validator verdict applied to that shared initial candidate; the guarded outcome equals the final result of the guarded pipeline after deterministic validation and any permitted repair. The preregistration also declared that independent constrained‑prompt arms were not executed; therefore any within‑pair differences can arise only from deterministic validator actions and permitted repair calls, not from different first‑pass samples or different initial prompts.

Model configuration, sampling notation, and budgetary constraints. The declared archived model is gpt‑5‑mini (provider recorded as openai\_responses) and the adapter enforced a strict call budget: planned\_episode\_count = 400, initial\_generation\_calls\_per\_task = 1, maximum\_repair\_calls\_per\_task = 1, maximum\_model\_calls = 400. The archived execution model\_configuration.json records temperature = null (unset). We explicitly state: temperature = null/unset is not evidence of deterministic hosted‑model sampling and does not imply temperature = 0. Because provider sampling semantics are external to the adapter, we treat the single shared initial generation per intent as the executed stochastic sample; subsequent guarded repairs (if any) are additional provider calls permitted only by the predeclared repair policy.

Deterministic verifier design and scoring rule. The deterministic\_netops\_validator\_v5 is a rule‑based verifier implemented to mechanistically evaluate intent\ensuremath{\rightarrow}configuration candidates. Its functional components (all coded as deterministic checks applied to the candidate configuration and the emulated Mininet state) include:
1) syntactic parsing and command‑pattern conformance (token‑level grammar and per‑vendor command templates);
2) required‑command sequencing and workflow policy checks driven by the per‑task payload (e.g., ordered steps, protected interfaces, minimum‑active constraints);
3) reachability and topology consistency checks that compare the candidate's intended final state against the emulated network topology and device states;
4) ACL and policy encodings derived from task payload constraints (checks for forbidden field changes or prohibited sequences that would violate policy). The verifier returns a structured trace: validation\_before (diagnostic list and violation codes), a binary validity verdict (valid: true/false), and, when violations are present and repair is permitted by policy, a localized diagnostic summary that is included verbatim in the repair prompt. The scoring rule full\_intent\_constraint\_validation yields success = 1 only when the final validated configuration satisfies all mechanizable constraints.

Repair policy and prompt semantics. The guarded pipeline applies a conservative bounded repair policy as preregistered: if and only if the verifier reports violations on the shared initial candidate, at most one verifier‑triggered repair call may be issued. Repair prompts include (a) the original shared initial candidate, (b) the verifier's localized diagnostic messages (violation codes and minimal context identifying implicated commands/fields), and (c) an explicit instruction to produce a corrected configuration that satisfies the listed constraints. Repairs were allowed only when the verifier flagged violations and were restricted to a single additional hosted‑model call per intent to preserve the predeclared budget and to evaluate cost per avoided violation.

Scoring, exclusions, and missingness treatment. The primary confirmatory analysis treats the pairwise binary success indicators (no detected violations after the pipeline completes for baseline and guarded) per the preregistration. Predeclared exclusion rules removed duplicate tasks detected via deterministic prompt+context hashing and episode pairs aborted due to unrecoverable infrastructure/container substitution failures; such exclusions would be reported (none occurred in the archived run). Ambiguous verifier outputs (diagnostics flagged by deterministic ambiguity rules) were conservatively treated as violations in the primary analysis (this conservative rule was preregistered). The missingness plan treated genuine infrastructure/API failures as missing‑at‑random for the primary test; a sensitivity analysis treating failures as conservative failures was predeclared.

Analysis plan and resampling specifics. The preregistered primary inferential procedure is the exact McNemar two‑sided test applied to the 2\ensuremath{\times}2 paired contingency table (n11, n10, n01, n00). Secondary and exploratory estimators declared in the preregistration include: (i) a paired bootstrap percentile 95\% confidence interval for the paired risk difference (bootstrap\_seed = 1729; bootstrap\_resamples = 10,000), and (ii) descriptive cost metrics (mean extra hosted model calls per intent and mean extra calls per avoided violation). We note that bootstrap percentile intervals can be unstable when discordant pair counts (n10 + n01) are small; this instability motivated treating the McNemar exact test as the formal confirmatory inference and labeling the bootstrap interval exploratory in the manuscript.

Provider accounting, auditability, and reproducibility artifacts. The experiment's adapter recorded a provider‑call audit: planned maximum\_model\_calls = 400; hosted\_model\_call\_event\_count was auditable in the execution manifest (the adapter records stage\_counts and condition\_counts). The preregistration required recording all raw provider responses, verifier traces, repair prompts, deterministic task manifests, and an execution manifest; these artifacts were archived for machine verification. Analysis code, seeds, and exact resampling parameters were archived to enable deterministic replication of the reported numerical analyses.

Why shared\_initial\_candidate semantics matters for causal interpretation. We emphasize that because the archived run used shared\_initial\_candidate semantics and permitted at most one repair call per intent, the paired comparison isolates the effect of deterministic verification and the bounded repair step on the shared initial sample. In this design baseline is not an independent generation arm; it is the validator's direct verdict on the sampled candidate. Any between‑condition improvements within a pair therefore must arise from validator‑triggered repair converting a failing shared initial candidate into a passing final configuration, not from differing initial samples or alternative prompting.

Implementation and operational controls. The adapter enforced deterministic task ordering and recorded task indices prior to model calls. Repair prompts were templated and deterministic given the verifier diagnostics; repair calls were logged with the exact repair prompt text and provider response. The deterministic verifier and task generator versions used in the run (deterministic\_netops\_validator\_v5 and deterministic\_netops\_task\_generator\_v5) are recorded in the archived model and artifact manifests to permit exact replay of the pipeline on the same synthetic task bank and emulation state. All numerical quantities reported in Results (model\_call counts, stage\_counts, and contingency table cells) are drawn from those archived manifests and reconciled in the deterministic\_reconciliation artifact.

\section{Results}
Execution completion and deterministic accounting. The preregistered run completed without infrastructure exclusions: planned\_episode\_count = 400, completed\_episode\_count = 400, complete\_pair\_count = 200, failed\_episode\_count = 0 (execution manifest). Audited provider accounting (deterministic\_reconciliation) reports hosted\_model\_call\_event\_count = 204; this decomposes deterministically into stage\_counts = \{shared\_initial\_generation: 200, repair: 4\} and condition\_counts recorded as \{shared: 200, guarded: 4\} in the manifest. Cache statistics show cache\_hit\_count = 0 and cache\_miss\_count = 204, confirming that every required hosted call executed and that the four repair calls were distinct provider calls. All deterministic consistency checks passed in the reconciliation artifacts.

Paired binary outcomes and primary confirmatory test. The preregistered paired contingency (analysis/paired\_contingency\_table.csv) is n11 = 196, n10 = 4, n01 = 0, n00 = 0. These cells yield baseline\_success\_rate = 196/200 = 0.98 and guarded\_success\_rate = 200/200 = 1.00; the paired risk difference (guarded - baseline) = 0.02. The preregistered exact McNemar two‑sided test on these paired indicators returned p = 0.125 (analysis results). Under the preregistered inferential contract this p‑value means we do not reject the null at \ensuremath{\alpha} = 0.05.

Repair events, candidate reuse, and per‑intent cost accounting. The stage decomposition (200 shared initial generation events + 4 repair events) matches the model‑call audit and demonstrates that the guarded pipeline invoked repairs for exactly four intents. From these audited counts we compute total\_model\_calls = 204, mean\_model\_calls\_per\_intent = 204/200 = 1.02, mean\_extra\_calls\_per\_intent\_due\_to\_repairs = 4/200 = 0.02. Because each repair converted one failing baseline into a passing guarded outcome in the archived contingency, the observed mean\_extra\_calls\_per\_avoided\_violation = 4 repairs / 4 avoided violations = 1.0. These arithmetic derivations follow directly from the archived execution manifest and deterministic reconciliation artifacts.

Shared\_initial\_candidate evidence and representative artifact substantiation. Representative task artifacts included in the evidence bundle show explicit reuse of the same sampled initial generation by both baseline and guarded episodes. For example, the representative tasks (task‑000004, task‑000005, task‑000006, etc.) exhibit identical shared\_initial\_candidate contents and identical shared\_initial response file SHA‑256 values across their baseline and guarded response artifacts. The archived validator\_trace fields for the discordant pairs record validation\_before entries showing violations and validation\_after entries showing repair\_applied = true with final validation\_after.valid = true; the bundle therefore substantiates that the four discordant cases were resolved by verifier‑triggered repairs operating on the identical initial candidate.

Diagnostic pattern observed in discordant pairs. In every archived discordant pair the validator\_trace shows a mechanizable violation detected in validation\_before (violation\_count > 0 and violation\_codes returned by deterministic\_netops\_validator\_v5) and a subsequent repair application that produced a final configuration satisfying full\_intent\_constraint\_validation. The manifest and per‑task scoring artifacts indicate that repairs were only issued for intents where the verifier produced explicit, localized diagnostics; no repair calls occurred for intents without detected violations. The archived traces therefore support the causal pathway: identical initial candidate \ensuremath{\rightarrow} verifier detects mechanizable violation \ensuremath{\rightarrow} single permitted repair call \ensuremath{\rightarrow} final guarded configuration passes deterministic validation.

Exploratory bootstrap interval and interpretation caution. The archived exploratory paired bootstrap (bootstrap\_seed = 1729; resamples = 10,000) produced a 95\% percentile bootstrap CI for the paired risk difference = [0.005, 0.04]. We report this interval as exploratory only. Practically, the bootstrap percentile CI in paired binary settings relies on resampling paired outcomes and can be unstable when the total number of discordant pairs (n10 + n01) is small; here n10 + n01 = 4 is small, which reduces the bootstrap's reliability as a standalone uncertainty summary. Accordingly, we treat the exact McNemar test as the prespecified confirmatory inference (p = 0.125) and use the bootstrap interval only to convey plausible effect‑size magnitudes suggested by the archived data.

Power context and observed baseline prevalence. The preregistered power calculation targeted a 10 percentage‑point absolute reduction with N \ensuremath{\approx} 157 (McNemar approximation) and noted that with N = 200 the detectable absolute risk difference under the same assumptions would be \ensuremath{\approx}8 percentage points. The observed baseline failure prevalence (2\%) was substantially lower than scenarios that motivated the target effect size, which reduced the study's power to detect small absolute improvements and produced the small discordant count (4) observed. This lower than anticipated baseline error rate is an empirical fact recorded in the archived condition summary and underlies the conservative interpretation of the null result from the preregistered test.

Audit reconciliation and reproducibility pointers. All numerical claims above (paired cell counts, model\_call totals, decomposition into shared initial and repair stages, bootstrap parameters, and seeds) are directly traceable to archived artifacts: execution/raw\_results.jsonl, analysis/paired\_contingency\_table.csv, deterministic\_reconciliation.json, execution/model\_configuration.json, and analysis/analysis\_log.jsonl. The execution manifest records started\_at\_utc and completed\_at\_utc timestamps for the run and the preregistration SHA‑256; the analysis artifacts record the exact bootstrap\_seed = 1729 and bootstrap\_resamples = 10000 used to produce the exploratory interval. Reproducibility of the primary contingency counts does not require re‑sampling because shared\_initial\_candidate semantics ensured exactly one initial sampled generation per intent was reused by both conditions; the archived response files and validator traces substantiate this reuse.

Summary of result‑section evidence. In short: (1) the formal preregistered confirmatory test did not reject the null hypothesis at \ensuremath{\alpha} = 0.05 (exact McNemar p = 0.125); (2) deterministic verifier‑triggered repairs converted four failing shared initial candidates into passing final configurations, as shown by per‑task validator\_trace entries; (3) the total hosted model‑call overhead for these repairs was low (4 additional calls across 200 intents, mean extra calls per intent = 0.02, mean extra calls per avoided violation = 1.0); (4) the exploratory bootstrap CI [0.005, 0.04] suggests small plausible effect magnitudes but must be interpreted cautiously given the sparse discordant pair count. These findings are reported verbatim from the archived execution and analysis bundle and are discussed further in the Discussion section.

\section{Discussion}
Causal interpretation under shared\_initial\_candidate semantics. Because execution used shared\_initial\_candidate semantics (one sampled initial generation per intent reused by baseline and guarded), paired improvements arise solely from deterministic validator/repair actions. The preregistration explicitly declared that independent constrained‑prompt arms were not executed; thus the observed n10 = 4 discordant pairs reflect validator‑triggered repairs that converted failing shared initial candidates into passing final configurations. Any statements attributing improvements to prompt‑engineering or differing first‑pass samples would be incorrect for this run.

Construct validity and verifier coverage. The deterministic\_netops\_validator\_v5 covers syntactic correctness, required command sequencing, reachability/topology checks against the emulated state, and encoded ACL/policy checks derived from per‑task payload constraints. This verifier can only detect violations expressible in its rules; higher‑order semantic misinterpretations not captured by those encodings remain undetectable and are identified in our limitations. Increasing verifier expressivity (formal policy ontologies, broader semantic checks) is a direct next step to reduce residual false negatives.

Internal validity and nondeterminism. The archived model\_configuration.json records temperature = null/unset; we reiterate that null/unset does not imply temperature = 0 or deterministic provider sampling. The execution manifest's provider audit (hosted\_model\_call\_event\_count = 204, stage\_counts, and condition\_counts) was reconciled deterministically (all deterministic consistency checks passed) and confirms that cross‑condition cache reuse did not occur; this audit supports the claim that baseline outcomes are the validator verdict on the identical shared initial generation used by guarded episodes.

Operational implications (bounded). Within the constrained Mininet emulation, synthetic templates, and validator rule set, a small repair budget (four repairs) removed the mechanically detectable baseline violations at very low hosted‑model cost. This artifact‑grounded observation suggests that when a deterministic verifier exists and provides localized actionable diagnostics, a bounded repair policy can remove template‑driven mechanically checkable faults in similar emulated settings. We do not claim production readiness, multi‑model generality, nor coverage for semantic errors outside the verifier.

Threats to validity and recommendations. Internal validity---shared\_initial\_candidate semantics were predeclared and executed; this prevents attributing gains to independent re‑sampling. Construct validity---the verifier cannot detect all semantic errors; extend verifier coverage and measure false negatives. External validity---single declared model (gpt‑5‑mini) and synthetic templates in Mininet limit generalization; evaluate additional hosted models and hardware testbeds. Conclusion validity---small discordant counts reduce power; future preregistered work could target higher baseline failure prevalence scenarios or increase N. Based on archived evidence we recommend: (1) expanding deterministic verifier expressivity to formalize more semantic/policy ontologies; (2) performing paired designs that also execute independent multi‑sample initial generations to separate prompt/sampling effects from repair effects; and (3) repeating provenance‑traced evaluations across multiple hosted models and template families.

\section{Limitations}
\begin{itemize}
\item Scope limited to Mininet‑style emulation with synthetic, provenance‑traced intent\ensuremath{\rightarrow}configuration tasks; results do not generalize directly to production hardware or live operations.
\item Single declared model (gpt‑5‑mini) evaluated; multi‑model generality is not established.
\item Deterministic validator coverage limited to mechanically checkable errors (syntax, reachability, ACL/policy encodings); higher‑level semantic or specification misinterpretations outside the validator may remain undetected.
\item Shared\_initial\_candidate semantics imply observed paired differences derive only from verifier‑triggered repair; this design does not measure effects of alternative prompting strategies or independent multi‑sample candidate selection.
\item Observed confirmatory effect (2 percentage points) did not reach statistical significance at \ensuremath{\alpha} = 0.05 for N = 200 (exact McNemar p = 0.125); low baseline failure prevalence (2\%) reduced power to detect small absolute improvements under some preregistered scenarios.
\end{itemize}

\section{Conclusion}
We preregistered and executed a provenance‑traced paired confirmatory study (C1‑GVR‑multivendor‑2026‑v1) using shared\_initial\_candidate semantics to measure whether a deterministic verifier plus bounded repair reduces mechanically detectable policy violations for intent\ensuremath{\rightarrow}configuration tasks in a Mininet emulation. The preregistered exact McNemar two‑sided test did not reject at \ensuremath{\alpha} = 0.05 (p = 0.125). Guarded GVR invoked four repairs and corrected the four observed baseline violations at low model‑call overhead (model calls = 204; mean extra calls per intent = 0.02), yielding a paired risk difference estimate of 0.02 and an exploratory paired bootstrap CI = [0.005, 0.04] (bootstrap\_seed = 1729, resamples = 10,000). These artifact‑grounded results indicate that when a deterministic verifier can express and check relevant constraints and provide localized feedback, a small repair budget can eliminate mechanically detectable policy violations in similar template‑based emulated settings. The results do not establish production readiness or multi‑model generality. All reported numbers, provider accounting, validator traces, and analysis artifacts are drawn verbatim from the archived execution and analysis bundles referenced in the preregistration and execution manifests.

References
[1] R. Mondal, A. Tang, R. Beckett, T. Millstein, and G. Varghese, "What do LLMs need to Synthesize Correct Router Configurations?", Proceedings of the ACM, 2023, doi: 10.1145/3626111.3628194.
[2] R. Cadenas, "Convergent Correctness in Stochastic Code Generation: A Generate‑Verify‑Repair Architecture for Deterministic Validation of Autonomous AI Coding Agents in Regulated Environments", SSRN, 2026, doi: 10.2139/ssrn.6754899.
[3] Z. Ding, "Executable but Wrong: Verifier‑Grounded Contracts and Counterexamples for LLM‑Generated Music Programs", 2026, doi: 10.31224/7995.
[4] C. Wang, M. Scazzariello, A. Farshin, S. Ferlin, D. Kostić, and M. Chiesa, "NetConfEval: Can LLMs Facilitate Network Configuration?", Proceedings of the ACM, 2024, doi: 10.1145/3656296.
[5] S. Y. Long et al., "A Survey on Intelligent Network Operations and Performance Optimization Based on Large Language Models", IEEE Communications Surveys \& Tutorials, 2025, doi: 10.1109/comst.2025.3526606.
[6] L. Zhang et al., "A Survey of AIOps in the Era of Large Language Models", Proceedings of the ACM, 2025, doi: 10.1145/3746635.

\section*{Disclosure Statement}
Declared model (archived): gpt‑5‑mini (provider recorded as openai\_responses). Adapter family: hosted\_netops\_gvr\_v1. Deterministic validator: deterministic\_netops\_validator\_v5. Task generator: deterministic\_netops\_task\_generator\_v5. Preregistration: C1‑GVR‑multivendor‑2026‑v1 (SHA‑256 = 9e0fc992\allowbreak{}66583597\allowbreak{}5b01289f\allowbreak{}7841926e\allowbreak{}f6abf358\allowbreak{}4252c5d0\allowbreak{}459ba0ec\allowbreak{}f166bb50), recorded prior to any model calls. Initial master prompt archived at provenance/master\_prompt.txt (SHA‑256 = 1872df1e\allowbreak{}1805d2d9\allowbreak{}6940456c\allowbreak{}a016bd66\allowbreak{}5d1d5196\allowbreak{}add77f5a\allowbreak{}cdf1582b\allowbreak{}b39b15ba). Human role: a human Research Director selected the topic family and constrained the initial master prompt prior to the autonomy lock; after the lock the AI autonomously performed literature verification, preregistration, experiment execution, analysis, peer review, revision, and manuscript drafting; no human manually edited technical manuscript content after the autonomy lock. Sampling semantics: execution used shared\_initial\_candidate (exactly one sampled initial generation per intent was produced and reused by both baseline and guarded); archived model configuration records temperature = null/unset (null/unset is not evidence of deterministic provider sampling and does not imply temperature = 0). Provider accounting (audited): hosted\_model\_call\_event\_count = 204 decomposed as 200 shared initial generations + 4 repair calls. Reported numbers, provider‑call accounting, validator traces, and analysis artifacts are drawn verbatim from the archived execution and analysis bundles referenced in the preregistration and execution manifests.

\begin{thebibliography}{99}
\bibitem{ref1} Rajdeep Mondal, Alan Tang, Ryan Beckett, Todd Millstein, George Varghese, "What do LLMs need to Synthesize Correct Router Configurations?", 2023, 10.1145/3626111.3628194
\bibitem{ref2} Rafael Cadenas, "Convergent Correctness in Stochastic Code Generation: A Generate-Verify-Repair Architecture for Deterministic Validation of Autonomous AI Coding Agents in Regulated Environments", 2026, 10.2139/ssrn.6754899
\bibitem{ref3} Zhengyang Ding, "Executable but Wrong: Verifier-Grounded Contracts and Counterexamples for LLM-Generated Music Programs", 2026, 10.31224/7995
\bibitem{ref4} Changjie Wang, Mariano Scazzariello, Alireza Farshin, Simone Ferlin, Dejan Kostić, Marco Chiesa, "NetConfEval: Can LLMs Facilitate Network Configuration?", 2024, 10.1145/3656296
\bibitem{ref5} S.Y. Long, Jingjing Tan, Bomin Mao, Fengxiao Tang, Yangfan Li, Ming Zhao, Nei Kato, "A Survey on Intelligent Network Operations and Performance Optimization Based on Large Language Models", 2025, 10.1109/comst.2025.3526606
\bibitem{ref6} Lingzhe Zhang, Tong Jia, Mengxi Jia, Yifan Wu, Aiwei Liu, Yong Yang, Zhonghai Wu, Xuming Hu, Philip S. Yu, Ying Li, "A Survey of AIOps in the Era of Large Language Models", 2025, 10.1145/3746635
\end{thebibliography}

\end{document}
